\section{Not Lost, but Sent Before}

Before his own death a Christian sees others die, and the Church teaches him how to grieve for them. The Apostle forbids one kind of sorrow: ``And we will not have you ignorant brethren, concerning them that are asleep, that you be not sorrowful, even as others who have no hope. For if we believe that Jesus died and rose again: even so them who have slept through Jesus, will God bring with him.''\endnote{1 Thessalonians 4:12--13, Douay--Rheims (Challoner), Vulgate numbering.}

Ambrose spoke of this grief at his brother's bier, before his own people. He had wept, and he did not repent of it. ``Not all weeping proceeds from unbelief or weakness. Natural grief is one thing, distrustful sadness is another, and there is a very great difference between longing for what you have lost and lamenting that you have lost it. Not only grief has tears, joy also has tears of its own. \ldots{} I confess, then, that I too wept, but the Lord also wept. He wept for one not related to Him, I for my brother. He wept for all in weeping for one, I will weep for thee in all, my brother.''\endnote{Ambrose, \emph{De excessu fratris sui Satyri} I.10, trans. H. de Romestin, in \emph{Nicene and Post-Nicene Fathers}, second series, vol.~10.}

\subsection{Cyprian in the plague}

The other sorrow, the distrustful sadness that treats the dead as though they were no more, is the one Cyprian of Carthage fought. He wrote \emph{On the Mortality} while a pestilence was carrying off Christians and pagans alike, and he would not let his people mourn their dead as lost. He tells them that God had often charged him to declare it ``that our brethren who are freed from this world by the Lord's summons are not to be lamented, since we know that they are not lost, but sent before; that, departing from us, they precede us as travellers, as navigators are accustomed to do; that they should be desired, but not bewailed; that the black garments should not be taken upon us here, when they have already taken upon them white raiment there.''\endnote{Cyprian, \emph{De mortalitate} 20, trans. R. E. Wallis, \emph{Ante-Nicene Fathers}, vol.~5.}

For Cyprian death is a passing over: ``That in the meantime we die, we are passing over to immortality by death; nor can eternal life follow, unless it should befall us to depart from this life. That is not an ending, but a transit, and, this journey of time being traversed, a passage to eternity. Who would not hasten to better things?''\endnote{\emph{De mortalitate} 22, Wallis.} He tells his people to look toward the home they are going to: ``We should consider, dearly beloved brethren---we should ever and anon reflect that we have renounced the world, and are in the meantime living here as guests and strangers. Let us greet the day which assigns each of us to his own home, which snatches us hence, and sets us free from the snares of the world, and restores us to paradise and the kingdom.''\endnote{\emph{De mortalitate} 26, Wallis.}

Others are waiting there: ``There a great number of our dear ones is awaiting us, and a dense crowd of parents, brothers, children, is longing for us, already assured of their own safety, and still solicitous for our salvation.''\endnote{\emph{De mortalitate} 26, Wallis.} The Church teaches that the saints care in this way for those still on earth. The Council of Trent bids the pastors teach that the saints, reigning together with Christ, offer their prayers to God for men, and that it is good and useful to call upon them.\endnote{Council of Trent, session XXV (1563), decree on the invocation, veneration, and relics of saints, DS 1821: \latin{Sanctos, una cum Christo regnantes, orationes suas pro hominibus Deo offerre}.} The Catechism sets beside the teaching of the Second Vatican Council two sayings of saints who knew they were dying. Saint Dominic told his brethren: ``Do not weep, for I shall be more useful to you after my death and I shall help you then more effectively than during my life.'' Saint Th\'er\`ese said: ``I want to spend my heaven in doing good on earth.''\endnote{CCC 956, citing \emph{Lumen gentium} 49, Saint Dominic's words to his brethren as he was dying, and Th\'er\`ese of Lisieux, \emph{Final Conversations}.} The communion of saints continues past the grave: the blessed in heaven are, in the words Paul VI made the Church's profession, ``interceding for us and helping our weakness by their fraternal concern.''\endnote{Paul VI, \emph{Credo of the People of God} 29, as quoted in CCC 1053.} The meeting with them, and the joy of it, belong to \emph{Then Face to Face}.
